\chapter{Extending the CloudSC Loop-Nest Analysis}
\label{ch:cloudsc}
% RQ4. Builds on Bonsall's semester project "Evaluating the Auto-Vectorization Amenability
% of DaCe-Generated Code" (References/DaCeCodegenAnalysis_AlexanderBonsall.pdf).

\section{Prior Results}
% Five loop nests: autoconversion snow, ice supersaturation adjustment, LU solver,
% rain evaporation (Abel-Boutle), saturation calculation.
% Variants: original Fortran, DaCe Fortran frontend, DaCe Python frontend.
% Findings: row/column-major porting bug in the Python frontend (LU solver 8.4x);
% forcing vectorization rarely helps; recommendations: stage reductions through scalar
% temporaries, FP flags for GCC or fewer connectors.

\section{What Was Left Open}
% 1. Warm-up not accounted for: variants still drifting at the end of the run.
% 2. DaCe timer resolution is 1 us; ice supersaturation (~16--19 us) cannot be ordered.
% 3. Kernels only C++ vectorized (DaCe failed) not all analysed.
% 4. Why -fassociative-math changes GCC's behaviour with no FP change in the assembly.
% 5. Rule "FP flags for GCC" not applied to rain evaporation (min/max-heavy).
% 6. Heavily branching loops (ice supersaturation, autoconversion) not studied.
% 7. Recommendations stated but not implemented in the code generator.

\section{Extended Setup}
% Reuse the LLR-40 protocol: ns timer outside DaCe, warm-up until steady state,
% min-of-k, NumPy/Fortran oracle. Optionally add numba/C/C++ representations of the loop
% nests to connect to RQ1.

\section{Results}
% One subsection per open item addressed.

\section{Implications for DaCe Code Generation}
